% ===== §8 Conclusions and Envoi =====
\section{Conclusions and Envoi}
\label{p18:sec:conclusions}

\noindent
This paper has isolated an injustice the theory of justice had not named: a wrong done not to what a subject has, or is taken to be, or can make known, or could come to be, but to her being, in the precise sense of her continuing activity as a generative subject. In the manner of the series, it closes not on a single synthesis but on the several registers in which the finding can be stated---each true, none reducible to the others, their plurality not a hedge but the form a finding takes when it sits at the meeting of an ontology, a formalism, a deontic structure, and an ethic.

\subsection{The findings, in their registers}
\label{p18:sub:conc-registers}

\emph{Ontologically}, the finding is that there is a mode of injustice consisting in the deprivation of being-as-generation. Because, for a relational subject, to be is to go on generating, the foreclosure of its generation is not the loss of one of its activities but the hollowing of its existence; and this can be done to a subject already fully formed, which is what distinguishes it in register from the constitutive injustice it descends from. The theory of justice acquires, with this layer, a floor beneath the floor it had supposed final---and the descent that reaches it is licensed not by analogy but by the identity of being with generation that the series argues on independent grounds.

\emph{Formally}, the finding is that effacement is a phase transition with an exact threshold. A subject loses its own frequency when the coupling between two unlike rhythms exceeds a bound fixed by their difference; the depth of the loss is set by the asymmetry of force; and---the result that organises all the rest---the loss occurs even under perfectly equal power. The good of intimacy is thereby shown to lie not in the strength of the bond but in whether it leaves each subject turning at its own frequency, and the long-standing identification of closeness with strength of coupling is exposed as the precise error that lets devotion efface.

\emph{Deontically}, the finding is that the justice answering this wrong is not additive but subtractive. It is not the transfer of any share, recognition, or hearing, but a negative and structural obligation to refrain from determining the domain in which another generates---an obligation whose content is non-determination and whose discharge is the leaving of room. This is the rarest shape a duty of justice can take: not \emph{give}, but \emph{withhold}; not the completion of the other's world, but the deliberate leaving-incomplete of one's own occupation of it.

\emph{Ethically}, the finding is that this obligation is owed in virtue of the other's standing as a generative being, on an authority that is hers; that it is detected not by rule but by a cultivated perception, since it is breached through devotion and announces itself by no transgression; and that it forbids, on pain of self-refutation, its own paternalistic enforcement. The ethic is thus doubly disciplined---deontological in ground and virtue-theoretic in application, and reflexively barred from becoming the management of the other it would protect.

\begin{claim}[The closing thesis]
Between two generative beings, justice requires not the transfer of value but the non-foreclosure of the other's capacity to generate. A giving that occupies the whole of the shared space is not the fulfilment of love but the quiet extinction of one of its two subjects; and the love that would keep both subjects in being must therefore learn a giving that withholds---must leave, within the field it generates, the room in which the other generates too. Justice, on this account, is not what one does for the other but what one leaves undone, that she may do it herself.
\end{claim}

\subsection*{Envoi}
\noindent
Two rhythms, unlike in their frequencies, may be brought so close that they beat as one---and in the beating-as-one the slower is not joined but silenced, kept in motion by a measure no longer its own. The bond looks then most like union at the very moment it has become most solitary, one rhythm sounding through two bodies. The whole of this paper has been an argument that the cure for that solitude is not to slow the faster rhythm but to loosen the bond between them, until each, turning again at its own measure, can be near the other without being overtaken by it---close enough that something is generated between them, free enough that it is generated by them both. That is the only nearness in which two generative beings remain two. It is, this paper has argued, the form that justice takes when what is at stake is no longer what we have, nor how we are seen, nor whether we can be heard, but whether, beside each other, we continue to be.
